% GENERATED FILE. Edit canonical lessons, then run npm run generate:books.
% canonical-source-hash: fnv1a64:0cd6b1f14164aceb
% canonical-lessons: TA-C423-nakaratci, TA-C423-puranakar, TA-C423-kutiyiruppu, TA-C423-vattaram, TA-C423-maiyam, TA-R423-review-416-423

\chapter{Parts of a City}
\label{ch:ta-parts-of-a-city}
\hlchaptermodality{423}

\begin{chapteropening}
\textbf{By the end of this chapter:} \emph{I can talk about a municipality, suburbs, housing colonies, localities and the centre of a town.}

Complete the last lesson of chapter 423: I can talk about a municipality, suburbs, housing colonies, localities and the centre of a town.
\end{chapteropening}

\section[\texorpdfstring{\ta{நகராட்சி} (nakarāṭci)}{நகராட்சி (nakarāṭci)}]{\ta{நகராட்சி} (nakarāṭci) — a municipality}

\label{lesson:TA-C423-nakaratci}



\begin{quote}
Before the new one: say the Tamil for binoculars, then the Tamil for to untie.
\end{quote}

\subsection*{You'll want to know: \ta{நகராட்சி}}
\textbf{\ta{நகராட்சி}} — \emph{nakarāṭci} — \textquotedblleft{}a municipality\textquotedblright{}.

\begin{quote}
\ta{நகராட்சி} \ta{புதிய} \ta{பூங்காவைத்} \ta{திறந்தது}.
\end{quote}

\begin{quote}
\emph{nakarāṭci putiya pūṅkāvait tiṟantatu.}
\end{quote}

\begin{quote}
The municipality opened a new park.
\end{quote}

\ta{நகர்}, town, + \ta{ஆட்சி}, rule.

\begin{cousinweb}
\textbf{\ta{நகராட்சி}} (\emph{nakarāṭci}) is built from \ta{நகர்}, \textquotedblleft{}a town\textquotedblright{} (DEDR 3568), and \ta{ஆட்சி}, \textquotedblleft{}rule, government\textquotedblright{} (DEDR 5157).
\end{cousinweb}

\begin{grammarlens}[title={What you've built}]
One more word for getting around.
\end{grammarlens}

\subsection*{Your turn}
\begin{itemize}
\raggedright
  \item \emph{Say it:} \emph{nakarāṭci}
  \item \emph{Say it:} \emph{nakarāṭci}, once more
  \item \emph{Say it:} \emph{nakarāṭci putiya pūṅkāvait tiṟantatu.} — The municipality opened a new park.
  \item \emph{Recall:} say the Tamil for binoculars, then the Tamil for to untie, then say \emph{nakarāṭci} again
\end{itemize}

\begin{tcolorbox}[breakable,colback=teal!4,colframe=teal!35!black,title={Before you move on}]
What does \ta{நகராட்சி} mean? (\textquotedblleft{}A municipality\textquotedblright{}.) Say it once more.
\end{tcolorbox}
\section[\texorpdfstring{\ta{புறநகர்} (puṟanakar)}{புறநகர் (puṟanakar)}]{\ta{புறநகர்} (puṟanakar) — a suburb}

\label{lesson:TA-C423-puranakar}



\begin{quote}
Before the new one: say the Tamil for to untie, then the Tamil for a municipality.
\end{quote}

\subsection*{You'll want to know: \ta{புறநகர்}}
\textbf{\ta{புறநகர்}} — \emph{puṟanakar} — \textquotedblleft{}a suburb\textquotedblright{}.

\begin{quote}
\ta{என்} \ta{அண்ணன்} \ta{புறநகரில்} \ta{வாழ்கிறான்}.
\end{quote}

\begin{quote}
\emph{eṉ aṇṇaṉ puṟanakaril vāḻkiṟāṉ.}
\end{quote}

\begin{quote}
My elder brother lives in a suburb.
\end{quote}

\ta{புறம்}, the outside, + \ta{நகர்}, town.

\begin{cousinweb}
\textbf{\ta{புறநகர்}} (\emph{puṟanakar}) is built from \ta{புறம்}, \textquotedblleft{}the outside\textquotedblright{} (DEDR 4333), and \ta{நகர்}, \textquotedblleft{}a town\textquotedblright{} (DEDR 3568).
\end{cousinweb}

\begin{grammarlens}[title={What you've built}]
One more word for getting around.
\end{grammarlens}

\subsection*{Your turn}
\begin{itemize}
\raggedright
  \item \emph{Say it:} \emph{puṟanakar}
  \item \emph{Say it:} \emph{puṟanakar}, once more
  \item \emph{Say it:} \emph{eṉ aṇṇaṉ puṟanakaril vāḻkiṟāṉ.} — My elder brother lives in a suburb.
  \item \emph{Recall:} say the Tamil for to untie, then the Tamil for a municipality, then say \emph{puṟanakar} again
\end{itemize}

\begin{tcolorbox}[breakable,colback=teal!4,colframe=teal!35!black,title={Before you move on}]
What does \ta{புறநகர்} mean? (\textquotedblleft{}A suburb\textquotedblright{}.) Say it once more.
\end{tcolorbox}
\section[\texorpdfstring{\ta{குடியிருப்பு} (kuṭiyiruppu)}{குடியிருப்பு (kuṭiyiruppu)}]{\ta{குடியிருப்பு} (kuṭiyiruppu) — a housing colony; a residence}

\label{lesson:TA-C423-kutiyiruppu}



\begin{quote}
Before the new one: say the Tamil for a municipality, then the Tamil for a suburb.
\end{quote}

\subsection*{You'll want to know: \ta{குடியிருப்பு}}
\textbf{\ta{குடியிருப்பு}} — \emph{kuṭiyiruppu} — \textquotedblleft{}a housing colony; a residence\textquotedblright{}.

\begin{quote}
\ta{புதிய} \ta{குடியிருப்பில்} \ta{ஒரு} \ta{பள்ளி} \ta{இருக்கிறது}.
\end{quote}

\begin{quote}
\emph{putiya kuṭiyiruppil oru paḷḷi irukkiṟatu.}
\end{quote}

\begin{quote}
There is a school in the new housing colony.
\end{quote}

\ta{குடி}, a home, + a noun made from \ta{இரு}, to stay.

\begin{cousinweb}
\textbf{\ta{குடியிருப்பு}} (\emph{kuṭiyiruppu}) is built from \ta{குடி}, \textquotedblleft{}a house, a home\textquotedblright{} (DEDR 1655), and the noun of \ta{இரு}, \textquotedblleft{}to be, to stay\textquotedblright{} (DEDR 480).
\end{cousinweb}

\begin{grammarlens}[title={What you've built}]
One more word for getting around.
\end{grammarlens}

\subsection*{Your turn}
\begin{itemize}
\raggedright
  \item \emph{Say it:} \emph{kuṭiyiruppu}
  \item \emph{Say it:} \emph{kuṭiyiruppu}, once more
  \item \emph{Say it:} \emph{putiya kuṭiyiruppil oru paḷḷi irukkiṟatu.} — There is a school in the new housing colony.
  \item \emph{Recall:} say the Tamil for a municipality, then the Tamil for a suburb, then say \emph{kuṭiyiruppu} again
\end{itemize}

\begin{tcolorbox}[breakable,colback=teal!4,colframe=teal!35!black,title={Before you move on}]
What does \ta{குடியிருப்பு} mean? (\textquotedblleft{}A housing colony; a residence\textquotedblright{}.) Say it once more.
\end{tcolorbox}
\section[\texorpdfstring{\ta{வட்டாரம்} (vaṭṭāram)}{வட்டாரம் (vaṭṭāram)}]{\ta{வட்டாரம்} (vaṭṭāram) — a locality; a district}

\label{lesson:TA-C423-vattaram}



\begin{quote}
Before the new one: say the Tamil for a suburb, then the Tamil for a housing colony.
\end{quote}

\subsection*{You'll want to know: \ta{வட்டாரம்}}
\textbf{\ta{வட்டாரம்}} — \emph{vaṭṭāram} — \textquotedblleft{}a locality; a district\textquotedblright{}.

\begin{quote}
\ta{இந்த} \ta{வட்டாரத்தில்} \ta{நிறைய} \ta{தோட்டங்கள்} \ta{இருக்கின்றன}.
\end{quote}

\begin{quote}
\emph{inta vaṭṭārattil niṟaiya tōṭṭaṅkaḷ irukkiṉṟaṉa.}
\end{quote}

\begin{quote}
There are many gardens in this locality.
\end{quote}

\ta{வட்டார} \ta{மொழி} is a local dialect.

\begin{cousinweb}
\textbf{\ta{வட்டாரம்}} (\emph{vaṭṭāram}), \textquotedblleft{}a locality, a circle of villages\textquotedblright{}, is a Tamil word whose origin is not settled.
\end{cousinweb}

\begin{grammarlens}[title={What you've built}]
One more word for getting around.
\end{grammarlens}

\subsection*{Your turn}
\begin{itemize}
\raggedright
  \item \emph{Say it:} \emph{vaṭṭāram}
  \item \emph{Say it:} \emph{vaṭṭāram}, once more
  \item \emph{Say it:} \emph{inta vaṭṭārattil niṟaiya tōṭṭaṅkaḷ irukkiṉṟaṉa.} — There are many gardens in this locality.
  \item \emph{Recall:} say the Tamil for a suburb, then the Tamil for a housing colony, then say \emph{vaṭṭāram} again
\end{itemize}

\begin{tcolorbox}[breakable,colback=teal!4,colframe=teal!35!black,title={Before you move on}]
What does \ta{வட்டாரம்} mean? (\textquotedblleft{}A locality; a district\textquotedblright{}.) Say it once more.
\end{tcolorbox}
\section[\texorpdfstring{\ta{மையம்} (maiyam)}{மையம் (maiyam)}]{\ta{மையம்} (maiyam) — a centre}

\label{lesson:TA-C423-maiyam}



\begin{quote}
Before the new one: say the Tamil for a housing colony, then the Tamil for a locality.
\end{quote}

\subsection*{You'll want to know: \ta{மையம்}}
\textbf{\ta{மையம்}} — \emph{maiyam} — \textquotedblleft{}a centre\textquotedblright{}.

\begin{quote}
\ta{நகரத்தின்} \ta{மையத்தில்} \ta{ஒரு} \ta{பெரிய} \ta{சந்தை} \ta{இருக்கிறது}.
\end{quote}

\begin{quote}
\emph{nakarattiṉ maiyattil oru periya cantai irukkiṟatu.}
\end{quote}

\begin{quote}
There is a big market in the centre of the city.
\end{quote}

Also an institution: \ta{பயிற்சி} \ta{மையம்}, a training centre.

\begin{cousinweb}
\textbf{\ta{மையம்}} (\emph{maiyam}) is usually traced to Sanskrit \emph{madhya}, \textquotedblleft{}the middle\textquotedblright{} (CDIAL 9804).
\end{cousinweb}

\begin{grammarlens}[title={What you've built}]
One more word for getting around.
\end{grammarlens}

\subsection*{Your turn}
\begin{itemize}
\raggedright
  \item \emph{Say it:} \emph{maiyam}
  \item \emph{Say it:} \emph{maiyam}, once more
  \item \emph{Say it:} \emph{nakarattiṉ maiyattil oru periya cantai irukkiṟatu.} — There is a big market in the centre of the city.
  \item \emph{Recall:} say the Tamil for a housing colony, then the Tamil for a locality, then say \emph{maiyam} again
\end{itemize}

\begin{tcolorbox}[breakable,colback=teal!4,colframe=teal!35!black,title={Before you move on}]
What does \ta{மையம்} mean? (\textquotedblleft{}A centre\textquotedblright{}.) Say it once more.
\end{tcolorbox}
\section[(dialogue)]{Looking back: from Boats and Oars to Parts of a City}

\label{lesson:TA-R423-review-416-423}



\begin{quote}
Say the Tamil for a boatman and an oar.
\end{quote}

\subsection*{Your turn}
Say these aloud:

\begin{itemize}
\raggedright
  \item a boatman, an oar, a dinghy, a catamaran, a net
  \item a sailor, an anchor, a raft, to flounder, a shore
  \item a lodge, a receptionist, an elevator, bedding, to spread out
  \item an air conditioner, hot water, to wrap up, to relax, hospitality
  \item laundry, a service, a porter, a wage, cleaning
  \item a guide, a handbook, to show around, to visit, to see off
  \item a camp, a torch, a flask, binoculars, to untie
  \item a municipality, a suburb, a housing colony, a locality, a centre
\end{itemize}

\begin{tcolorbox}[breakable,colback=teal!4,colframe=teal!35!black,title={Before you move on}]
A boatman? (\textbf{\ta{படகோட்டி}}, \emph{paṭakōṭṭi}.) An oar? (\textbf{\ta{துடுப்பு}}, \emph{tuṭuppu}.) A dinghy? (\textbf{\ta{தோணி}}, \emph{tōṇi}.) A catamaran? (\textbf{\ta{கட்டுமரம்}}, \emph{kaṭṭumaram}.) A net? (\textbf{\ta{வலை}}, \emph{valai}.) A sailor? (\textbf{\ta{மாலுமி}}, \emph{mālumi}.) An anchor? (\textbf{\ta{நங்கூரம்}}, \emph{naṅkūram}.) A raft? (\textbf{\ta{தெப்பம்}}, \emph{teppam}.) To flounder? (\textbf{\ta{தத்தளி}}, \emph{tattaḷi}.) A shore? (\textbf{\ta{கரை}}, \emph{karai}.) A lodge? (\textbf{\ta{விடுதி}}, \emph{viṭuti}.) A receptionist? (\textbf{\ta{வரவேற்பாளர்}}, \emph{varavēṟpāḷar}.) An elevator? (\textbf{\ta{மின்தூக்கி}}, \emph{miṉtūkki}.) Bedding? (\textbf{\ta{படுக்கை}}, \emph{paṭukkai}.) To spread out? (\textbf{\ta{விரி}}, \emph{viri}.) An air conditioner? (\textbf{\ta{குளிரூட்டி}}, \emph{kuḷirūṭṭi}.) Hot water? (\textbf{\ta{வெந்நீர்}}, \emph{vennīr}.) To wrap up? (\textbf{\ta{போர்த்து}}, \emph{pōrttu}.) To relax? (\textbf{\ta{இளைப்பாறு}}, \emph{iḷaippāṟu}.) Hospitality? (\textbf{\ta{விருந்தோம்பல்}}, \emph{viruntōmpal}.) Laundry? (\textbf{\ta{சலவை}}, \emph{calavai}.) A service? (\textbf{\ta{சேவை}}, \emph{cēvai}.) A porter? (\textbf{\ta{சுமைதூக்கி}}, \emph{cumaitūkki}.) A wage? (\textbf{\ta{கூலி}}, \emph{kūli}.) Cleaning? (\textbf{\ta{துப்புரவு}}, \emph{tuppuravu}.) A guide? (\textbf{\ta{வழிகாட்டி}}, \emph{vaḻikāṭṭi}.) A handbook? (\textbf{\ta{கையேடு}}, \emph{kaiyēṭu}.) To show around? (\textbf{\ta{சுற்றிக்காட்டு}}, \emph{cuṟṟikkāṭṭu}.) To visit? (\textbf{\ta{பார்வையிடு}}, \emph{pārvaiyiṭu}.) To see off? (\textbf{\ta{வழியனுப்பு}}, \emph{vaḻiyaṉuppu}.) A camp? (\textbf{\ta{முகாம்}}, \emph{mukām}.) A torch? (\textbf{\ta{டார்ச்}}, \emph{ṭārc}.) A flask? (\textbf{\ta{குடுவை}}, \emph{kuṭuvai}.) Binoculars? (\textbf{\ta{தொலைநோக்கி}}, \emph{tolainōkki}.) To untie? (\textbf{\ta{அவிழ்}}, \emph{aviḻ}.) A municipality? (\textbf{\ta{நகராட்சி}}, \emph{nakarāṭci}.) A suburb? (\textbf{\ta{புறநகர்}}, \emph{puṟanakar}.) A housing colony? (\textbf{\ta{குடியிருப்பு}}, \emph{kuṭiyiruppu}.) A locality? (\textbf{\ta{வட்டாரம்}}, \emph{vaṭṭāram}.) A centre? (\textbf{\ta{மையம்}}, \emph{maiyam}.)
\end{tcolorbox}
